\scanpage{5}
mod $p$ des valeurspropres de l'endomorphisme donné de $E(Y)$ ). Enfin, (et c'est un petit luxe supplémentaire moins important pour nous, car seul les trois derniers termes de (8.1.3) sont pour nous vraiment importants du point de vue assymptotique ...) \deleted{\unsure{on aura} $\mathrm{Ker}\,u$} le \deleted{\unsure{premeir}}premier terme $\mathrm{Ker}\,u_0$ dans (8.1.3) est nul:
\[
\mathrm{Ker}\,u_0=0\;.
\tag*{(8.3.3)}
\]

8.4. Notons maintenant que l'on a des applications semi-linéaires $\mathrm{Frob}$ opérant fonctoriellement sur tous les espaces de cohomologie envisagés, et en particulier sur \deleted{\unsure{ceux}} \inserted{les termes} intervenant dans (8.1.3), qui sont donc des vectoriels de dimension finie sur $k$ avec endomorphisme $p$-linéaire. Considérons, pour un tel vectoriel $V$, le plus grand sous-espace vectoriel stable $V_{ss}$ (partie ``semi-simple'' de $V$) stable par $F_V$, et sur lequel $F_V$ soit \inserted{``semi-simple'' i.e.} surjectif (NB $k$ est parfait); \deleted{Le foncteur} il revient au,même de dire que $F$ est surjectif sur ce sous-espace, et \deleted{\unsure{nul sur le}}\uline{nilpotent} sur le quotient. Pour $V$ variable, la formation de $V_{ss}$ est un foncteur exact en $V$\unsure{.} Ainsi, la suite exacte (8.1.3) nous donne, compte tenu de (8.3.3), la suite exacte d'espace vectoriels à opérateurs $p$-linéaires semi-simples:
\[
\begin{aligned}
0&\handwritten{\longrightarrow}(\mathrm{Ker}\,u_1)_{ss}\handwritten{\longrightarrow}(\mathrm{Ker}\,u_2)_{ss}
\inserted{\handwritten{=({}_p(\mathrm{Tors}^n(X)))_{ss}}}\handwritten{\longrightarrow}\\
& (E(Y)\otimes_W k)_{ss}\handwritten{\longrightarrow}E_{\rm DR}(Y)_{ss}\handwritten{\longrightarrow}
\overset{\handwritten{\substack{0\\\Vert}}}{(\mathrm{Coker}\,u_2)_{ss}}\handwritten{\longrightarrow}0\;.
\end{aligned}
\tag*{(8.4.1)}
\]
\margin{\handwritten{en vertu de 8.4.4.}}

\inserted{\handwritten{8.4.2.}} D'après la théorie de structure des vectoriels à opération $p$-linéaire semi-simple, \deleted{qui} \inserted{ceux-ci} peuvent s'interpréter aussi comme des Modules libres de type fini sur le faisceau d'anneaux étale constant $\underline{\mathrm F}_p$ (corps premier à $p$ éléments), ou encore \deleted{\unsure{comme}} (une fois choisie une clôture séparable) comme des \deleted{\unsure{$G$-$\underline{\mathrm F}$}} représentations de $G=\mathrm{Gal}(\bar k/k)$ sur des vectoriels de dimension finie sur $\handwritten{\unsure{\mathrm F|\mathbb F|\underline{\mathrm F}}_p}$. Dans cette catégorie il y a une autodualité, qui s'interprète indffféremment en termes de dualité d'espaces de dimension finie sur $k$ avec opération $p$-linéaire semi-simple, ou en termes de dualité de
% archive sheet number: 4
